\section{From the formulas to an implemented comparison}
\label{sec:implemented-study}
The earlier studies exposed two gaps. The branch counter missed the work of reading leaf
masks, and the short-key recall experiment did not measure a native key index. The next study
implements all three methods with the same score, then checks the pieces before comparing
complete requests. This section reports that study separately from the historical results.

\subsection{What is shared, and what changes}
A, B and C use the same C++ lookup-table scorer and the same top-$K$ heap. Each query remains
floating point; each document stores one sign bit per coordinate. The final score is the full
weighted binary score from Section~\ref{sec:foundations}, with smaller document IDs first when
scores tie. There is no float reranker in this comparison. Thus recall means recovery of the
exact binary-score top 100, not relevance to a person reading the documents.

The Python driver supplies cached query vectors and selected cluster IDs. A fresh process
builds one configuration, releases its input mappings, warms up a query, then measures routing
and search together. The native search also records its own elapsed time. Measuring the
outer call includes the Python/native boundary; reporting only the native timer would omit it.

The machine is an Apple M5 Max with 128 GiB of physical memory. Each worker uses one CPU
thread and handles one query at a time. The comparison budget is 32 GB, in decimal bytes.
The machine's physical capacity and the experiment's memory constraint are different numbers.
Power configuration is recorded before and after each worker. Equal settings do not establish
fixed CPU frequency or a completely idle machine, so timings are repeated in separate processes.

\subsection{The implemented key directory}
The native C build sorts rows by key inside each cluster and stores their full codes in that
same order. Its existing document IDs are the postings. There is no extra $N$-element array of
posting references. This is the layout alternative described after Equation~\eqref{eq:keys-storage}.

The directory is a C++ \texttt{unordered\_map}. It uses separate chaining: occupied entries are
allocated as nodes and reached through a bucket-pointer array. Equation~\eqref{eq:hash-storage}
describes a different, open-addressed layout and must not be used as its measured allocation.
Let $L_j$ be the number of hash buckets, $b_{\rm node}$ the allocated node size including fields
and alignment, and $b_{\rm ptr}$ a bucket-pointer size. Then the implemented allocation is
\begin{align}
M_{C,\rm allocated}={}&M_{\rm codes,capacity}+M_{\rm IDs,capacity}
 +\sum_j\left(U_j b_{\rm node}+L_j b_{\rm ptr}\right)\notag\\
 &+M_{\rm containers}+M_{\rm allocator}.
\label{eq:native-key-storage}
\end{align}
The last two terms cover container objects and allocator bookkeeping. They depend on the
standard library and build. The logical directory payload is easier to check exactly: each
occupied entry holds a key and a start/count pair,
\[
M_{\rm directory,payload}=\sum_j U_j\left(b_k+2b_o\right).
\]
In this build those fields occupy 4, 8 and 8 bytes respectively. Their 20-byte sum is a payload
count, not a claim that an allocated hash node occupies only 20 bytes.

B retains packed document rows and a second, transposed bitplane representation. A is built
without bitplanes. Giving A the unused plane allocation would distort the memory comparison.
All three retain the IDs needed to return original document rows.

\subsection{The implemented key search uses one stream}
The implementation has the same free key coordinates in every selected cluster. It enumerates
one weighted sequence of keys and looks up each key in every selected cluster. This differs
from the per-cluster target example earlier in the report.

\begin{lstlisting}[caption={The native key-search policy used in the measured study.}]
SearchKeys(query, selected_clusters, K, candidate_target, lookup_limit):
    build the shared full-score lookup table
    order the key coordinates by increasing absolute query weight
    queue = [the ideal key, with mismatch penalty zero]
    best = empty top-K heap
    while queue is not empty:
        if best is full and the best remaining score bound cannot improve it:
            stop
        state = remove the lowest-penalty flip set
        key = ideal key XOR state.flip_mask
        for cluster in selected_clusters:
            if the lookup limit is reached: return best
            count one lookup, including an empty lookup
            for document in directory[cluster][key]:
                score all document coordinates; offer the score and ID to best
        if a nonzero candidate target is reached: return best
        add the next flip sets to the queue
    return best in score/ID order
\end{lstlisting}
The score bound is the full ideal score minus twice the key's mismatch penalty. Coordinates
outside the key are allowed their best possible contribution. This remains a safe upper bound
when routing fixes some of those coordinates, although it may be loose. The implementation
does not yet tighten it separately for each cluster.

The candidate target is checked after finishing a key in all selected clusters. A lookup limit
can stop in the middle of that cluster loop. Both limits are approximations; zero disables a
limit. The exact score bound remains available in either case. Equal-score candidates require
the shared ID rule, so the bound includes conservative floating-point slack.

There is one enumeration queue, not $P$ separate queues. Charge its preparation and generated
states once; charge each cluster lookup separately. Every returned candidate still receives all
$G$ score-table terms. The tail-only and full-key score reuse discussed earlier are possible
alternatives, but they are not credited to the measured implementation.

\subsection{Check the work before fitting milliseconds}
Logical code, ID, plane and key-directory bytes are checked against separate Python formulas.
Native array capacity is reported separately because a vector may reserve more entries than it
currently holds. The tests also compare returned IDs against an independent full-score sort.
The reference calculation starts from document signs, not from the native packed scorer.

B reports split-word visits and leaf-word visits separately. C reports generated keys, directory
lookups including misses, and fully scored documents. Both report their largest tracked queue
or mask allocation. These counts explain a timing result; a low document-score count alone does
not establish that the whole query performed little work.

\begin{examplebox}{Two cuts do not shorten the physical bitmap}
Use every possible eight-bit document once, so $N=256$ and $W=\ceil{256/64}=4$ words.
The first three query coordinates have magnitude one; the remaining five have magnitude
$0.0001$. Request four results and stop splitting when 64 documents remain.

Two matching cuts leave $256/2^2=64$ documents. The search visits the root, its preferred child,
and the leaf: three nodes. It reads $2\times4=8$ split words and another four words to enumerate
the leaf. It then fully scores 64 documents. At the second split, one queued alternative,
the parent and both new children occupy four masks, or $4\times4\times8=128$ bytes.

The native measurement matches each count and returns the independently calculated top four.
With a one-node budget it returns no candidates: that neighboring truncated case checks that
the test can observe missing results. This example establishes the stated path, not a universal
bound of two cuts for arbitrary query weights.
\end{examplebox}

\subsection{Fit elapsed time, then try different sizes}
Instruction counts cannot be multiplied by a universal ``nanoseconds per operation.'' A lookup
that hits a nearby cache line and a lookup that fetches a new memory page have different costs.
The first timing model therefore uses a small set of measured work features with nonnegative
coefficients. Its search-time inputs are:
\begin{itemize}
\item A: a fixed term and the number of full score-table terms.
\item B: the same terms, split plus leaf word visits, and visited nodes.
\item C: the same score terms, directory lookups, and a queue-work estimate.
\end{itemize}
The coefficients combine effects of arithmetic, memory access and result selection. They are
not individual CPU-instruction timings. Measured routing time is added per query before taking
a median; summing the medians of separate stages would not give the median complete request.

There were 117 calibration configurations on random sign data with one cluster. An initial
fit minimized squared errors in milliseconds. Long queries dominated that objective and some
short branch queries had large relative errors. Refitting the same calibration data using
squared percentage errors improved the diagnostic comparison. The original fit remains saved.

The revised coefficients were then frozen and checked on 52 configurations at different sizes
and a new seed. All four scan settings and all 24 key settings had median prediction error
within 20\%. For B, 21 of 24 settings met that criterion; its worst error was 27.2\%.
The failures remain in the tables. These coefficients describe a limited range of working-set
sizes, not an established billion-document latency model.

This model also uses \emph{observed} work counts. Predicting how many branches or empty keys a
new query will visit is a different problem, governed by the document and query distribution.
A fit that converts measured work into milliseconds does not validate a forecast of recall or
candidate count. The later tuning comparisons use measured latency, rather than allowing a
poorly transferred model to choose a winner.

\subsection{Give the baseline enough parameter choices}
The real-data study reuses 256-bit Nomic/MS MARCO document codes at $N=10{,}000$, $100{,}000$
and $1{,}000{,}000$. All methods get the same candidate router families and cluster layouts;
each may select its own fastest tested configuration at a target recall.

A coarse 100,000-document run initially suggested a B advantage near 95\% recall. The scan
baseline's next available probe setting jumped from 92.76\% recall to a full scan. Adding
intermediate probe counts removed that apparent advantage. This was a gap in the experiment,
not evidence that B's traversal was faster.

The later probe choices use the actual ranks of the correct neighbors' clusters. Let
$r_{qi}$ be the rank of the cluster containing correct neighbor $i$ for tuning query $q$.
For $Q$ queries and top $K$, the exact local-scan recall after probing $P$ clusters is
\begin{equation}
\widehat{R}_{\rm route}(P)=\frac{1}{QK}
       \sum_{q=1}^{Q}\sum_{i=1}^{K}\ind[r_{qi}\le P].
\label{eq:empirical-route-rank}
\end{equation}
Sorting these $QK$ ranks gives the smallest empirical cutoff for any requested target.
The actual router checks the proposed cutoff; tied cluster scores are handled explicitly.
This closes a known probe-grid gap without claiming that recall is determined by $P/C$.

\subsection{Freeze settings, then use different queries}
The refined runs were combined into 6,821 distinct settings. Equivalent configurations pool
their timing observations. Random exploration uses a fixed set of seeds, averaged per query;
the fastest lucky seed is not selected as the method's result. The chosen parameters were
frozen before evaluating 200 query IDs excluded from the original 300-query collection.

Two selection rules were retained. The \emph{mean} rule picks the fastest measured setting whose
tuning recall reaches the target. The \emph{conservative} rule requires the lower bootstrap
percentile of tuning recall to reach it. This second rule leaves more margin, but it is still
a rule fitted to a finite sample, not a guarantee about all future queries.

There are 72 distinct frozen settings. Their predeclared seeds, three process blocks and three
timing repeats produced 258 workers and 154,800 query observations. All workers completed and
kept the same recorded power configuration. The largest process lifetime peak was 231,129,088
bytes, below the 32 GB constraint for these measured runs. It includes construction and runtime
memory, so it is more conservative than the logical index payload alone.

Only 7 of 36 mean-selected target points reached their requested recall on the new queries.
All 36 conservative points did. Re-encoding three old query texts with the fresh-query recipe
exactly reproduced their cached vectors; no parameters were changed in response to the misses.
A tuning average that just reaches 99\% is not a promise of 99\% on another finite query sample.

\begin{center}\small
\begin{tabular}{@{}rr rrr rrr@{}}\toprule
 & & \multicolumn{3}{c}{Median milliseconds} & \multicolumn{3}{c}{Achieved recall (\%)}\\
$N$ & Target & A & B & C & A & B & C\\\midrule
10,000 & 80\% & 0.0979 & 0.1052 & 0.1129 & 84.91 & 84.91 & 84.91\\
10,000 & 90\% & 0.1339 & 0.1474 & 0.1560 & 91.79 & 91.79 & 91.79\\
10,000 & 95\% & 0.1683 & 0.1818 & 0.1968 & 95.84 & 95.84 & 95.84\\
10,000 & 99\% & 0.2330 & 0.2542 & 0.2812 & 99.27 & 99.27 & 99.27\\
100,000 & 80\% & 0.2345 & 0.2454 & 0.2698 & 80.55 & 80.55 & 80.55\\
100,000 & 90\% & 0.5034 & 0.5293 & 0.5918 & 90.98 & 91.29 & 91.29\\
100,000 & 95\% & 0.7527 & 0.7976 & 0.9308 & 95.47 & 95.47 & 95.57\\
100,000 & 99\% & 1.6219 & 1.6779 & 1.8109 & 99.44 & 99.44 & 99.44\\
1,000,000 & 80\% & 0.6749 & 0.6985 & 0.7506 & 82.53 & 82.53 & 82.53\\
1,000,000 & 90\% & 1.6341 & 1.6856 & 1.8802 & 91.73 & 91.73 & 91.73\\
1,000,000 & 95\% & 3.3600 & 3.4295 & 3.7055 & 95.47 & 95.47 & 95.47\\
1,000,000 & 99\% & 9.5796 & 9.9645 & 11.0127 & 99.25 & 99.25 & 99.25\\
\bottomrule\end{tabular}\normalsize\end{center}

The table reports the conservative choices. Each method can choose different cluster counts,
probe counts and local settings. ``Target'' is the common minimum requirement; their achieved
recalls need not be numerically identical. Exact settings and all mean-rule target misses are
saved alongside the measurements.

Uncertainty is calculated by resampling queries and whole process blocks, using the same draws
for the paired methods. Timing repetitions are not treated as independent recall observations.
A supported alternative speedup requires both settings to meet the target and memory/power
checks, and the paired 95\% interval for $t_A/t_{\rm alternative}$ to lie above one. None of the
48 evaluated comparisons meets that rule. These intervals describe individual comparisons;
they are not simultaneous guarantees across every setting.

The measured conclusion is therefore narrow but useful: for this cached binary score, data,
single-thread implementation and tested parameter space, the independent evaluation did not
establish a B or C speedup over the tuned scan baseline. It does not exclude different query
weight distributions, other data layouts, hardware kernels or larger corpora.

\subsection{How to reproduce this chain}
Start with \path{experiments/README.md}. The small configuration checks components without a
dataset download. The isolated calibration and confirmation configurations check the timing
model. The frozen real selection lives in \path{results/frozen-real-2026-09-10/}; its independent
measurements and uncertainty are in \path{results/independent-evaluation-2026-09-10/}.
Each run stores its configuration, source and native hashes, environment, per-query observations
and process status. Raw-case ZIPs are verified against every original file before archival.

The implementation follows a short path through the repository:
\begin{center}\small
\begin{tabularx}{\linewidth}{@{}lY@{}}\toprule
Paper topic & Code to inspect\\\midrule
Shared binary score & \path{cpp/score.hpp}\\
Full scan and bitplane traversal & \path{cpp/index.cpp}\\
Weighted key enumeration & \path{cpp/key_index.cpp}\\
One complete measured query & \path{experiments/worker.py}\\
Independent storage/work formulas & \path{experiments/models.py}\\
Pooled tuning choices & \path{experiments/choices.py}\\
Frozen-query replay and comparison & \path{experiments/evaluate.py}, \path{experiments/evaluation_report.py}\\\bottomrule
\end{tabularx}\normalsize\end{center}
